\section{Section algebras of cotangent extensions}
\label{sec:elliptic-cotangent-extension-algebra}

Continue with $M=\mathbb P_{\mathrm{lines}}(\mathcal O_E\oplus A)$, where $E=\mathbb C/\Lambda$, $\Lambda=\mathbb Z+\mathbb Z\tau$, $\operatorname{Im}\tau>0$, and $\deg A=d>0$. Use the class generators of \cref{lem:split-elliptic-cotangent-cohomology}. Let $W$ be an extension of $\mathcal O_M$ by $\Omega_M^1$ with class $\rho\xi+Kf^*\eta$, where $\rho\ne0$, and put $V=W^*$. The homogeneous version of the charts in \cref{thm:split-elliptic-cotangent-torsor} adds the quotient coordinate $w$. On the elliptic cover they change by
\[
 v=1/t,\qquad u_\infty=u,\qquad
 y_\infty=-t^2y-\rho t w,\qquad w_\infty=w.
\]
For $\lambda=m\tau+n$, the lattice acts by
\[
 (z,t,u,y,w)\longmapsto
 (z+\lambda,b_\lambda t,u+m(Kw+q_0ty),y/b_\lambda,w),
 \qquad q_0=2\pi\mathrm{i}d.
\]
Here $b_\lambda$ is the normalized factor in \cref{lem:elliptic-degree-automorphy-normalization}. Fiber-homogeneous functions on $W$ represent symmetric tensors of $V$. Their ruling-coordinate algebra retains the quadratic relation of the affine torsor, while the elliptic monodromy controls which weights descend.

\begin{lemma}\label{lem:elliptic-extension-fiber-ring}
On $\mathbb C\times\mathbb P^1$, the graded algebra of fiber-homogeneous polynomial functions on the pulled-back $W$ is
\[
 \mathcal O(\mathbb C)
 [c_-,c_+,h,w,u]/\bigl(c_-c_+-h(h+\rho w)\bigr),
\]
with every displayed variable of degree one. In the finite ruling chart,
\[
 c_-=y,\qquad c_+=t(ty+\rho w),\qquad h=ty.
\]
The lattice fixes $h,w$, multiplies $c_-$ by $b_\lambda^{-1}$ and $c_+$ by $b_\lambda$, and sends $u$ to $u+m(q_0h+Kw)$.
\end{lemma}
\begin{proof}
Set $s_1=y$ and $s_2=ty+\rho w$ in the finite chart. Since $\rho\ne0$, these are invertible linear coordinates in place of $y,w$, with $w=(s_2-ts_1)/\rho$. In the infinite chart put
\[
 (s_{1,\infty},s_{2,\infty})
       =(-v y_\infty-\rho w,-y_\infty).
\]
The transition is $(s_{1,\infty},s_{2,\infty})=t(s_1,s_2)$. Thus the two coordinates are those of $\mathcal O_{\mathbb P^1}(-1)^{\oplus2}$, while $u$ is the coordinate of a trivial line.

A polynomial homogeneous of degree $j$ in $s_1,s_2$ has ruling coefficients of degree at most $j$ in $t$, by holomorphy in the infinite chart. Every admissible monomial is a product of $s_1,s_2,ts_1,ts_2$: distribute its $t$-exponent among its $j$ fiber factors. These four generators satisfy the determinant relation
\[
 s_1(ts_2)-(ts_1)s_2=0.
\]
There are no further relations. Indeed, the dimension in degree $j$ of the function space is $(j+1)^2$, by counting the $j+1$ fiber monomials and $j+1$ ruling coefficients. The polynomial algebra in four variables modulo this nonzero quadratic relation has the same dimension
\[
 \binom{j+3}{3}-\binom{j+1}{3}=(j+1)^2,
\]
where the second term is zero for $j<2$. The surjective graded map is consequently an isomorphism in every degree. Including $u$ adjoins a polynomial variable. The argument is uniform in $z$ and all coefficient functions are entire.

Identify the four generators as $c_-=s_1$, $h=ts_1$, $h+\rho w=s_2$, and $c_+=ts_2$. Their determinant relation becomes the stated quadratic relation. They are polynomial also in the infinite chart, where
\[
 c_-=-v^2y_\infty-\rho v w,\qquad
 c_+=-y_\infty,\qquad h=-v y_\infty-\rho w.
\]
The lattice formulas follow by substitution in the finite chart, proving the assertions.
\end{proof}

\begin{theorem}\label{thm:elliptic-extension-polynomial-ring}
There is a degree-preserving isomorphism of complex algebras
\[
 \bigoplus_{r\geq0}H^0(M,\operatorname{Sym}^rV)
 \simeq
 \left(\bigoplus_{r\geq0}H^0(M,\operatorname{Sym}^rT_M)\right)[w],
\]
where the variable $w$ on the right has degree one.
\end{theorem}
\begin{proof}
Give $c_-$ weight $-1$, $c_+$ weight $1$, and $h,w,u$ weight zero. These weights come from the ruling action and are preserved by the quadratic relation and lattice. Reducing products $c_-c_+$ shows that a weight $-k$ polynomial is uniquely $c_-^kF(z,h,w,u)$, and a weight $k$ polynomial is uniquely $c_+^kF(z,h,w,u)$, for $k\geq1$. Here $F$ has remaining degree $r-k$ in tensor degree $r$. Weight zero consists of $F(z,h,w,u)$.

For uniqueness, embed the relation algebra into
\[
 \mathcal O(\mathbb C)[h,w,u,c_-,c_-^{-1}],
 \qquad c_+\longmapsto h(h+\rho w)/c_-.
\]
The reduced expressions have distinct powers of $c_-$; their coefficients can vanish only if the original coefficients vanish, since $h(h+\rho w)$ is a nonzero polynomial. This proves both the weight decomposition and its directness.

Let $U_3$ be the coefficient bundle of linear polynomials in $h,w,u$, with lattice rule
\[
 F(z+\lambda,h,w,u)
       =F(z,h,w,u-m(q_0h+Kw)).
\]
The negative weight $-k$ block is $H^0(E,A^k\otimes\operatorname{Sym}^{r-k}U_3)$, and the positive weight $k$ block has $A^{-k}$ in place of $A^k$. This follows from the factors of $c_-$ and $c_+$ in the preceding lemma, using the inverse translation of $u$ when expressing invariance. In particular, the signs of the line twists agree with \cref{lem:split-elliptic-tensor-weights}.

Since $q_0\ne0$, put $H=h+(K/q_0)w$. Then the translation of $u$ is $u\mapsto u-mq_0H$. The coefficient bundle is consequently
\[
 U_3\simeq U\oplus\mathcal O_Ew,
\]
where the variable $h$ of the previously defined $U$ is replaced by $H$. Its symmetric powers are the direct sums of $w^j\operatorname{Sym}^{s-j}U$ for $0\leq j\leq s$. The negative-twist vanishing in \cref{lem:elliptic-unipotent-symmetric-cohomology} eliminates every positive weight. The same lemma identifies the zero-weight section algebra with $\mathbb C[H,w]$.

In each negative weight, expansion in powers of $w$ identifies the remaining spaces with the corresponding tangent blocks in \cref{lem:split-elliptic-tensor-weights}, with a polynomial variable $w$ adjoined. Send their formal Euler variable $h$ to $H$ and their leading factor $y^k$ to $c_-^k$. This gives bijections in every tensor degree and weight. Products between surviving blocks multiply their coefficients and add their powers of $c_-$; the eliminated positive weights contribute no products. The bijections therefore preserve multiplication and give the asserted graded algebra isomorphism.
\end{proof}

\begin{corollary}\label{cor:elliptic-extension-symmetric-dimensions}
For every integer $r\geq0$,
\[
 h^0(M,\operatorname{Sym}^rV)
       =r+1+d\binom{r+3}{4}.
\]
The tautological line on the quotient projective bundle $\mathbb P(V)$ is big of volume $d$, and the section algebra of $V$ is not finitely generated. These conclusions apply to every K\"ahler canonical extension on $M$.
\end{corollary}
\begin{proof}
The polynomial-algebra isomorphism and \cref{thm:split-elliptic-symmetric-dimensions} give
\[
 h^0(M,\operatorname{Sym}^rV)
 =\sum_{j=0}^r h^0(M,\operatorname{Sym}^jT_M)
 =r+1+d\sum_{j=0}^r\binom{j+2}{3}
 =r+1+d\binom{r+3}{4}.
\]
Here $\binom{2}{3}=0$. The surface $M$ is projective, and coherent GAGA makes its holomorphic bundle $V$ algebraic. Thus $\mathbb P(V)$ is a projective fourfold, and its tautological volume is the limit of $24h^0(M,\operatorname{Sym}^rV)/r^4$, equal to $d>0$. Thus the line is big. A finitely generated section algebra of $V$ would have a finitely generated quotient by $w$, contradicting \cref{thm:split-elliptic-tangent-ring-nonfinite}.

For a K\"ahler canonical extension $0\to\mathcal O_M\to V\to T_M\to0$, the dual cotangent extension has nonzero ruling-fiber component by \cref{thm:split-elliptic-cotangent-torsor}. Its class therefore has $\rho\ne0$ in the conventions used here. The same argument applies, proving the last assertion.
\end{proof}
